\section{引言}
\label{sec:introduction}

随着大语言模型（Large Language Model, LLM）推理与指令遵循能力的不断增强，以其为核心构建的智能体（Agent）正成为人工智能领域的研究与应用热点。借助工具调用（Tool Calling）机制，智能体能够突破模型自身在知识与交互能力上的限制，访问外部API、数据库、文件系统和运维平台等异构服务，完成数据分析、业务办理、系统运维等复杂任务。ReAct~\cite{yao2023ReAct}、Toolformer~\cite{schick2023Toolformer}等工作确立了“推理—行动”相结合的智能体范式，为工具调用能力的规模化应用奠定了基础。近年来，模型上下文协议（Model Context Protocol, MCP）等标准化工具调用协议的提出~\cite{anthropic2025MCP}，进一步统一了工具描述、能力发现与调用交互规范，显著降低了跨平台工具接入与复用的成本，推动智能体工具生态朝着开放化、规模化和组件化的方向快速演进。然而，随着大量第三方工具、外部数据源和后端业务系统被纳入智能体执行链路，基于预先配置和静态白名单的传统信任机制已难以适应开放工具生态的安全需求，工具调用过程中的运行时安全风险日益凸显。

工具调用在增强智能体行动能力的同时，也将大语言模型的安全边界从内容生成层面扩展到真实业务系统与基础设施层面。一次工具调用往往涉及用户、智能体宿主、工具服务和后端资源等多个主体，调用请求、工具元数据与执行结果需要在不同信任域之间持续流转，链路中任一环节的不可信输入都可能沿调用链传播，造成跨信任域的安全影响。具体而言，风险主要来自两个方向：其一，恶意或被攻陷的工具服务可能通过篡改工具描述、伪造返回结果、在外部内容中嵌入提示注入指令等方式，干扰智能体的任务规划与工具选择，进而导致系统指令被覆盖、上下文被污染、记忆被篡改以及敏感信息泄露；其二，被提示注入劫持的智能体、配置不当的客户端或恶意调用方，也可能生成超出用户真实意图或授权范围的工具请求，引发越权资源访问、危险命令执行、批量数据导出和业务状态篡改等高风险操作，直接威胁后端工具服务、核心数据与业务系统的安全稳定。已有研究证实了提示注入攻击在工具调用场景中的危害性与传播能力~\cite{liu2024PromptInjection}，但如何对跨组件、跨信任域调用链路中的请求与响应进行统一理解和运行时约束，仍是智能体规模化部署所面临的关键安全挑战。

尽管已有研究与工业实践持续关注大模型工具调用的安全防护，但现有方案仍存在三方面不足。首先，防护视角多局限于单一安全边界，缺乏对工具调用链路整体风险的系统建模：部署于智能体侧的工作（如AgentArmor~\cite{wang2025AgentArmor}、StruQ~\cite{chen2025StruQ}）主要检测工具返回内容中的提示注入与恶意指令，较少关注智能体发出的请求是否会对后端系统造成越权访问或危险操作；而沿袭传统Web与API安全思路的方案（如IsolateGPT~\cite{wu2025IsolateGPT}、ACE~\cite{li2026ACE}、SAGA~\cite{syros2026SAGA}）主要依赖身份认证、接口白名单与基于角色的访问控制，难以理解调用请求背后的任务意图与语义风险，尚未形成覆盖请求生成、工具执行与结果返回全过程的统一运行时防护机制。其次，防护手段多依赖规则匹配、字段校验与静态权限策略，缺乏对调用语义及上下文关系的深入理解：既难以识别经同义改写、编码混淆或参数拆分的提示注入，也难以判断单次调用本身合法、却偏离当前任务目标的语义越权行为（例如用户仅要求查询单条记录时却执行全表导出）；Progent~\cite{shi2025Progent}、Pro2Guard~\cite{wang2025Pro2Guard}、AgentSpec~\cite{wang2026AgentSpec}、DRIFT~\cite{li2025DRIFT}和A2AS~\cite{neelou2025A2AS}等工作虽在可编程权限、形式化策略与运行时自防御等方面提升了特定场景的安全控制能力，但由于缺少对调用主体、任务意图、目标资源、操作影响等异构信息的统一语义表示，各防护环节之间难以共享风险判定结果与策略知识。最后，现有防护系统通常与特定部署位置或智能体框架紧密耦合，智能体侧与工具服务侧的防护组件往往各自实现协议解析、风险识别与策略决策逻辑，能力难以复用，不仅推高了开发与运维成本，也容易导致不同信任边界上策略表达不一致、风险信息无法关联、防护覆盖出现盲区。因此，亟需一种能够适配不同部署点位、复用核心语义分析与策略决策能力的运行时安全机制。

针对上述问题，本文提出一种面向智能体工具调用链路的运行时安全网关。该网关以跨信任域的工具调用过程为保护对象，通过统一的语义表示与策略决策内核，对调用请求、工具元数据和执行结果进行持续的分析与约束。具体来说，网关首先从工具调用消息及其上下文中抽取调用主体、任务意图、操作类型、目标资源、资源敏感级别和潜在影响等关键语义要素，将协议各异、形式异构的调用数据转化为统一的结构化表示；在此基础上，策略决策引擎结合部署位置、主体权限、资源属性和上下文约束对调用行为进行风险判定，并视情况采取放行、拒绝、重写、降权、人工审批或隔离执行等处置措施。

在部署方式上，本文并未将智能体侧与工具服务侧设计为两套相互独立的防护系统，而是采用“统一内核、双侧部署”的架构：网关既可以部署在智能体宿主与外部工具之间，检测工具描述、工具返回内容及外部数据中隐藏的提示注入与上下文污染；也可以部署在工具服务入口，对智能体发出的调用请求进行语义级授权与风险约束，识别越权访问、危险操作、批量数据导出和业务状态篡改等行为。两类部署点位共享同一套语义解析、风险表示与策略决策能力，仅根据所处信任边界加载相应的检测规则与执行策略，从而在保持部署灵活性的同时，避免防护能力的割裂。考虑到工具服务直接连接核心数据与业务资源，本文重点研究工具服务入口处的请求风险识别与运行时控制，并通过智能体侧的部署实验验证统一防护内核的可迁移性与可扩展性。

本文的主要贡献如下：

\begin{itemize}
\item \textbf{工具调用链路威胁模型。}面向由用户、智能体宿主、外部工具服务和后端资源构成的跨信任域调用链路，本文系统分析了不可信工具内容与不安全的调用请求在链路中的产生、传播与影响过程，建立了覆盖工具元数据污染、响应提示注入、上下文操控、语义越权、危险操作和数据泄露等风险的威胁模型。该模型明确了攻击者能力、信任边界、攻击入口与安全目标，为工具调用运行时防护的设计与评估提供了统一基础。

\item \textbf{统一语义驱动的运行时安全网关。}本文设计并实现了一种面向智能体工具调用链路的运行时安全网关。网关通过对调用主体、任务意图、操作类型、目标资源、敏感等级和行为影响进行结构化建模，实现了对异构工具调用消息的统一语义分析；在此基础上，网关以共享的策略决策与执行内核支持智能体出口与工具服务入口两类部署形态，并提供放行、拦截、参数重写、人工审批、隔离执行等细粒度处置机制，缓解了现有方案防护视角割裂、语义理解不足和核心能力难以复用的问题。

\item \textbf{原型实现与系统化实验评估。}本文实现了安全网关原型系统，并构建了覆盖工具描述污染、响应提示注入、越权资源访问、危险命令执行、敏感数据导出和业务状态篡改等风险的测试集。实验结果表明，网关在工具服务侧风险检测任务上取得了\textbf{[XX.X\%]}的准确率、\textbf{[XX.X\%]}的召回率和\textbf{[XX.X\%]}的F1值，在智能体侧提示注入检测任务上取得了\textbf{[XX.X\%]}的检测率；与基于关键词匹配和静态权限控制的基线方法相比，F1值分别提升了\textbf{[XX.X]}和\textbf{[XX.X]}个百分点。同时，网关处理单次工具调用引入的平均额外延迟仅为\textbf{[XX ms]}，P95延迟为\textbf{[XX ms]}，表明该方法能够在较低运行时开销下有效识别并阻断典型的工具调用风险。

\end{itemize}